\section{Restricted shuffle and row/column-uniform operator pipelines and anti-diagonal untwisting}\label{sec:duality}

We state the scope of this section up front. In \Cref{sec:osr} we analyze Operator Schmidt Rank
($\OSR$) across the row/column tensor cut $F_{\mathrm{row}}^d \otimes_F (F_{\mathrm{col}}^d \otimes_F F^b)$
for linear pipelines built from Stone's shuffle and unshuffle \cite{Stone1971} together with
row-uniform operators $I_d \otimes B$ and column-uniform operators $A \otimes I_{db}$. This invariant
applies \emph{only} to that restricted operator pipeline: it \emph{does not} answer
\cite[Problem~9.5]{SeventeenthPaper2026} and \emph{does not} yield a lower bound for Turing machines.
Indeed, Problem~9.5 of \cite{SeventeenthPaper2026} asks for a functional that is subadditive over
machine steps, large for the swap at every scale separately, and \emph{``not raised at all scales by
a single linear-time operation''}, whereas within-row cyclic shear has $\OSR = d$
(\Cref{prop:shear-control}(a)) yet is computed in $O(L)$ steps on a $2$-work-tape Turing machine
(\Cref{prop:stone-tm}). Within the specified operator class, the pass bound $t \ge m$ of
\Cref{thm:osr-pass}(a) is attained with zero row/column-uniform stages by Stone's identity
$S^m = \tau$ \cite{Stone1971} (\lean{stoneShuffleGridLinear_iterate_eq_swap}).

\subsection{Operator Schmidt Rank on restricted shuffle pipelines and linear-time machines}\label{sec:osr}

Write $\lean{GridState}(F, d, b) := (\mathrm{Fin}\,d \to \mathrm{Fin}\,d \to \mathrm{Fin}\,b \to F)$,
separating the row index $r \in \mathrm{Fin}\,d$ from the within-row column-and-lane index
$(c, \ell) \in \mathrm{Fin}\,d \times \mathrm{Fin}\,b$. An $F$-linear operator $K$ on
$\lean{GridState}(F, d, b)$ has \emph{Operator Schmidt Rank at most $R$}, written $\OSR(K) \le R$
(\lean{HasOperatorSchmidtRankLe}; see \cite{Nielsen2003,Tyson2003}), if there exist $R$ row matrices
$A_i \in \Mat_{d\times d}(F)$ and $R$ within-row linear maps $B_i \in \End_F(F^d \otimes_F F^b)$ such
that
\[
  (K X)_{r,c,\ell}
  \;=\;
  \sum_{i=0}^{R-1} \sum_{r'=0}^{d-1} (A_i)_{r,r'}\,\bigl(B_i(X_{r'})\bigr)_{c,\ell},
\]
and we write $\OSR(K) = R$ (\lean{HasExactOperatorSchmidtRank}) when $\OSR(K) \le R$ and no
$R' < R$ suffices.

\begin{lemma}[Submultiplicativity and unit-rank operators; cf.\ {\cite{Nielsen2003,Tyson2003}}]\label{lem:osr-submult}
\formal{OperatorDualityProofs.lean}{operatorSchmidtRankSubmultiplicative,hasOperatorSchmidtRankLe_comp}
\formal{DualityAndSwapConsequences.lean}{isRowOrColUniform_rankLe_one}
Over any commutative semiring $F$ (and in particular any field):
\begin{enumerate}[label=(\alph*)]
\item If $\OSR(K_1) \le R_1$ and $\OSR(K_2) \le R_2$, then $\OSR(K_2 \circ K_1) \le R_2 R_1$.
\item Every row-uniform column operator $I_d \otimes B$ (\lean{IsRowUniformColOp}, $(KX)_r = B(X_r)$)
  and every column-uniform row operator $A \otimes I_{db}$ (\lean{IsColUniformRowOp},
  $(KX)_{r,c,\ell} = \sum_{r'} A_{r,r'} X_{r',c,\ell}$) satisfies $\OSR(K) \le 1$.
\end{enumerate}
\end{lemma}
\begin{outline}
Expanding $K_2 \circ K_1$ over $(i_2, i_1) \in \mathrm{Fin}\,R_2 \times \mathrm{Fin}\,R_1$ multiplies
the row matrices $A_{2,i_2} A_{1,i_1}$ and composes the within-row maps $B_{2,i_2} \circ B_{1,i_1}$
using $R_2 R_1$ terms; $I_d \otimes B$ and $A \otimes I_{db}$ are single elementary tensors.
\end{outline}

\begin{theorem}[Operator Schmidt Rank of transpose and Stone's shuffle; cf.\ {\cite{Nielsen2003,Tyson2003}}]\label{thm:osr-swap}
\formal{OperatorDualityProofs.lean}{operatorSchmidtRankSwap,operatorSchmidtRankStoneShuffle}
Let $F$ be any field and let $b \ge 1$.
\begin{enumerate}[label=(\alph*)]
\item For every $d \ge 1$, the matrix transpose $\tau = \lean{gridSwapLinear}(F, d, b)$ satisfies
  $\OSR(\tau) = d^2$.
\item On $d = 2^m$, Stone's $1$-bit cyclic shuffle $S = \lean{stoneShuffleGridLinear}(F, m, b)$ and
  unshuffle $S^{-1} = \lean{stoneUnshuffleGridLinear}(F, m, b)$ are mutual inverses with
  $\OSR(S) \le 4$ and $\OSR(S^{-1}) \le 4$ for all $m \ge 0$, and for every $m \ge 1$ ($d = 2^m \ge 2$),
  \[
    \OSR(S) \;=\; \OSR(S^{-1}) \;=\; 4.
  \]
\end{enumerate}
\end{theorem}
\begin{outline}
In~(a), decomposing $\tau$ by source/target row pairs uses $d^2$ terms, and evaluating any $R$-term
decomposition on row-localized lane-$0$ basis inputs produces $R$ matrices spanning $\Mat_{d\times d}(F)$,
forcing $R \ge d^2$. In~(b), for $m \ge 1$, Stone's shuffle sends $(r, c)$ to $(r', c')$ by moving
the top bit of $r$ into the low bit of $c'$ and the top bit of $c$ into the low bit of $r'$;
splitting on $(\text{low bit of } c',\, \text{low bit of } r') \in \{0,1\}^2$ gives a $4$-term
decomposition whose four nonzero quadrant blocks are linearly independent once $2^m \ge 2$ (and
symmetrically for $S^{-1}$).
\end{outline}

\begin{theorem}[Pass lower bounds and attainment for restricted shuffle and row/column-uniform pipelines]\label{thm:osr-pass}
\formal{OperatorDualityProofs.lean}{stoneShufflePipelinePassLowerBound}
\formal{DualityAndSwapConsequences.lean}{boundedOSRPipelinePassLowerBound,stoneShuffleGridLinear_iterate_eq_swap}
Let $F$ be any field and let $b \ge 1$.
\begin{enumerate}[label=(\alph*)]
\item On $d = 2^m$ ($m \ge 0$), if an operator list $K_1, \dots, K_n$ has composition
  \lean{composeGridOps} equal to $\tau$, where $t$ indices are equal to $S$ or $S^{-1}$
  and every remaining operator satisfies \lean{IsRowUniformColOp} or \lean{IsColUniformRowOp}
  (of the form $I_d \otimes B$ or $A \otimes I_{db}$), then $(2^m)^2 \le 4^t$ and $t \ge m$.
  Conversely, iterating $S$ $m$ times satisfies $S^m = \tau$ (\lean{stoneShuffleGridLinear_iterate_eq_swap},
  over any semiring and for any $m, b \ge 0$), attaining $t = m$.
\item More generally, for any $d \ge 1$ and $R \ge 1$, if a composed pipeline equals $\tau$ with $t$
  operators satisfying $\OSR \le R$ and all remaining operators satisfying $\OSR \le 1$, then
  $d^2 \le R^t$, and for $R \ge 2$ one has $t \ge \dfrac{2\log d}{\log R}$.
\end{enumerate}
\end{theorem}
\begin{outline}
By \Cref{lem:osr-submult} and \Cref{thm:osr-swap}(b), the composed pipeline has $\OSR \le 4^t$
(respectively $\OSR \le R^t$), whereas $\OSR(\tau) = d^2$ by \Cref{thm:osr-swap}(a); taking real
logarithms when $R \ge 2$ gives $2\log d \le t\log R$. For attainment in~(a), induction on
$0 \le k \le m$ shows that $(S^k X)_{r,c,\ell} = X_{(c \bmod 2^k)2^{m-k} + \lfloor r/2^k\rfloor,\,
(r \bmod 2^k)2^{m-k} + \lfloor c/2^k\rfloor,\,\ell}$, giving $S^m = \tau$ at $k = m$.
\end{outline}

In the shared library \texttt{lean-machine} (\leanname{LeanMachine.ShuffleShear}), we construct
explicit transition tables for one Stone shuffle pass, one Stone unshuffle pass, and within-row
cyclic shear, and transport them into \leanname{Transpose.TM} in \leanfile{ShuffleShearBridge.lean}:

\begin{proposition}[Verified $O(L)$-time multitape Turing machines for shuffle, unshuffle, and row shear]\label{prop:stone-tm}
\formal{ShuffleShearBridge.lean}{stoneShuffleTM_threeTape_linear,stoneUnshuffleTM_twoTape_linear,rowShearTM_twoTape_linear}
With $b = \mathrm{entryBits}(d)$ (\leanname{Transpose.entryBits}) and
$L = \mathrm{inputLength}(d) = d^2 b$ (\leanname{Transpose.inputLength}):
\begin{enumerate}[label=(\alph*)]
\item \lean{stoneShuffleTransposeTM} has $3$ work tapes ($k = 3$) and, for every $m \ge 1$ ($d = 2^m$)
  and every matrix $A$ (\leanname{Transpose.Matrix}), computes \lean{stoneShuffleTMMatrix} on $(m, A)$
  within $150 L$ steps.
\item \lean{stoneUnshuffleTransposeTM} has $2$ work tapes ($k = 2$) and, for every $m \ge 1$ ($d = 2^m$)
  and every matrix $A$ (\leanname{Transpose.Matrix}), computes \lean{stoneUnshuffleTMMatrix} on $(m, A)$
  within $150 L$ steps.
\item \lean{rowShearTransposeTM} has $2$ work tapes ($k = 2$) and, for every $d \ge 1$ and every matrix
  $A$ (\leanname{Transpose.Matrix}), computes within-row cyclic shear \lean{rowShearTMMatrix} on
  $(d, A)$ (sending entry $(r, c, \ell)$ to $A_{r, (c+r)\bmod d, \ell}$) within $200 L$ steps.
\end{enumerate}
\end{proposition}
\begin{outline}
Because $d$ and $b$ are not hardcoded in the finite control, the machines in
\leanname{LeanMachine.ShuffleShear} first determine unary counters for $b$ and $d$ (or $H = d^2/2$)
from the input length $L = d^2 b$ in $O(L)$ steps. The $3$-work-tape shuffle machine stores unary
counters for $b$ and $H$ on two work tapes, buffers the first $H$ words on a third work tape, and
interleaves $b$-bit words from the buffer and the second half of the input tape. The $2$-work-tape
unshuffle machine emits even-indexed $b$-bit words directly from the input tape while buffering
odd-indexed words on tape~$0$ using a unary $b$-counter on tape~$1$, then replays tape~$0$. The
$2$-work-tape row-shear machine keeps a save/emit row template on tape~$0$: for row~$r$, it buffers
the first $rb$ bits on tape~$1$, streams the remaining $(d - r)b$ bits directly from the input tape
to the output, and then emits the $rb$ buffered bits from tape~$1$ (advancing the save region on
tape~$0$ by $b$ cells for the next row). \lean{computesWithin_toTransposeTM} transports each proof
identically to
\leanname{Transpose.ComputesWithin}.
\end{outline}

Two controls clarify the boundary of $\OSR$ and of single-cell information measures:

\begin{proposition}[Within-row shear $\OSR = d$ and row-diagonal source-arity preservation]\label{prop:shear-control}
\formal{OperatorDualityProofs.lean}{control_RowShearAndRowDiagonalAlternation}
\formal{DualityAndSwapConsequences.lean}{rowDiagonal_SourceArityGaugeInvariance}
Let $F$ be any field, $d \ge 1$, and $b \ge 1$.
\begin{enumerate}[label=(\alph*)]
\item Within-row cyclic shear \lean{rowShearLinear} on $(F, d, b)$ lies in the row-diagonal subalgebra
  \lean{IsRowDiagonalOp} ($(KX)_r = B_r(X_r)$) yet has operator Schmidt rank~$d$.
\item For any $s \in \N$, the output-row source-arity bound $\le s$ (\lean{HasOutputRowSourceArityLe})
  is preserved under two-sided composition $D_1 \circ K \circ D_0$ with arbitrary row-diagonal
  operators $D_0, D_1$ (over any commutative semiring $F$), whereas $\OSR(\mathrm{id}) \le 1 < d =
  \OSR(\mathcal{S}_{d,b} \circ \mathrm{id})$ separates $\OSR$ from row-diagonal gauge invariance for
  all $d \ge 2$ (writing $\mathcal{S}_{d,b} := {}$\lean{rowShearLinear}$(F, d, b)$).
\item On $d = 2^m$ ($m \ge 0$), every alternating product $D_t U_t D_{t-1} \cdots U_1 D_0$
  (\lean{alternatingRowDiagProduct}) of $t + 1$ row-diagonal operators $D_i$ and $t$ shuffle or
  unshuffle passes $U_j \in \{S, S^{-1}\}$ has output-row source arity at most $2^t$; if it equals
  $\tau$, then $2^m \le 2^t$ and $t \ge m$.
\end{enumerate}
\end{proposition}
\begin{outline}
In~(a), \lean{rowShearLinear}${} = \sum_{r=0}^{d-1} E_{r,r} \otimes \sigma^r$ (where
$(\sigma v)_{c,\ell} = v_{c+1,\ell}$), and linear independence
of the $d$ cyclic column shifts $\sigma^0, \dots, \sigma^{d-1}$ forces $\OSR = d$. In~(b) and~(c),
row-diagonal operators act within each row without mixing source rows, while each $S$ or $S^{-1}$
combines at most two source rows into each output row; since row~$0$ of $\tau(X)$ depends on all
$2^m$ input rows when $b \ge 1$, equality with $\tau$ forces $2^m \le 2^t$.
\end{outline}

Over finite fields, \Cref{prop:shear-control}(c) is a
linear special case of \cite[Proposition~6.25(a)]{SeventeenthPaper2026} (proved in text), which bounds
arbitrary (not necessarily linear) pipelines over a $q$-symbol alphabet on $d = 2^m$ whose stages read
input rows inside affine cosets of $\F_2^m$: row-diagonal maps have coset dimension $0$ and
$S, S^{-1}$ have coset dimension $1$, so \cite[Proposition~6.25(a)]{SeventeenthPaper2026} bounds
transposed sets by $q^{(d^2 + 2^t d)/2}$ and subsumes the conclusion $2^m \le 2^t$ of
\Cref{prop:shear-control}(c).
By contrast, \Cref{thm:osr-pass}(a) allows dense
column-uniform linear row operators $A \otimes I_{db}$ (which read all $2^m$ rows and have coset
dimension $m$ in \cite[Proposition~6.25(a)]{SeventeenthPaper2026}) while excluding row-varying
column operators such as \lean{rowShearLinear}.

\begin{proposition}[Skew one-time-pad zero-mutual-information control]\label{prop:otp-control}
\formal{OperatorDualityProofs.lean}{control_SkewOneTimePadInformationTrap}
\formal{DualityAndSwapConsequences.lean}{uniformShannonMI_eq_zero_of_indep,skewOneTimePad_ShannonMutualInfoZero}
For $d \ge 1$ and $b \ge 1$, define $W(X) := X^{\T} \oplus X$ on $\{0,1\}^{d\times d\times b}$
(\lean{skewXorPad}). Then $W(X) \oplus X = X^{\T}$ for every $X$, and under uniformly random $X$,
every single entry $W_{r,c}$ is statistically independent of every source column $X_{:, c'}$, so its
discrete Shannon mutual information (\lean{uniformShannonMI}) satisfies
\[
  I\bigl(W_{r,c}\;;\; X_{:,c'}\bigr) \;=\; 0
  \qquad
  \text{for all } r, c, c' \in \Z_d.
\]
\end{proposition}
\begin{outline}
When $r = c$, $W_{r,r} = 0$ is constant. When $r \ne c$, the two entries $X_{c,r}$ and $X_{r,c}$ lie
in distinct columns $r \ne c$, so at least one lies outside column $c'$ and acts as an independent
uniform pad on the other, giving exact product-count independence and zero Shannon mutual
information.
\end{outline}

\subsection{Anti-diagonal untwisting of shift-collision equations}\label{sec:antidiag}

We close this section by recording a \leanname{TensorProduct} packaging of the shift-covariance
subspace of \cite[Lemma~3.1]{SharedDecoders2026} (and the shift-link covariance spaces of
\cite[Section~6]{SeventeenthPaper2026}) together with the compatibility characterization of
\cite{SharedDecoders2026}'s Lean library. In \cite{SharedDecoders2026}, Lemma~3.1
(\leanname{CapacityAmbient.lean}) establishes the necessity inclusion $C_E \subseteq H_E$
(\leanname{SharedDecoders.IsSharedCode.commonPreimage_le_covariance}), the anti-diagonal row
normalization (\leanname{SharedDecoders.covarianceSpace}, \leanname{SharedDecoders.covarianceEquiv}),
its dimension $bd\chi(E)$ (\leanname{SharedDecoders.finrank_covarianceSpace}), and its rotation
invariance (\leanname{SharedDecoders.covarianceSpace_rot}), while the two-way compatibility
equivalence (\leanname{SharedDecoders.isSharedCode_iff_compatible}) is proved in
\leanname{CapacityGluing.lean} of \cite{SharedDecoders2026}'s Lean library.
Theorem~\ref{thm:anti-diagonal}(a),(c)
and the upper bound of \Cref{prop:mask-floors}(a) below restate those properties; what is new here
is the packaging of the anti-diagonal coordinate change $a = c + j \in \Z_d$ (\lean{antiDiagLinear},
\lean{fromAntiDiagLinear}),
\[
  Z_{c,a} \;:=\; X_{c,\,a - c},
  \qquad
  X_{c,j} \;=\; Z_{c,\,j + c}
  \qquad (c, a, j \in \Z_d),
\]
for $d \ge 1$ as the explicit linear tensor isomorphism
$\Psi = \lean{antiDiagTensorEquiv} : (\Z_d \to F) \otimes_F (\Z_d \to A) \xrightarrow{\;\sim\;}
\Mat_{d\times d}(A)$ sending $u \otimes v$ to $(c, j) \mapsto u_c\,v_{j+c}$, together with the
converse construction of \Cref{prop:mask-floors}(a) and \Cref{prop:mask-floors}(b).

\begin{theorem}[Pairwise shift-collision characterization and anti-diagonal tensor separation; cf.\ {\cite[Lemma~3.1]{SharedDecoders2026}} and \leanname{SharedDecoders.isSharedCode_iff_compatible}]\label{thm:anti-diagonal}
\formal{OperatorDualityProofs.lean}{pairwiseShiftCollisionCharacterization,antiDiagonalTensorSeparation}
\formal{DualityAndSwapConsequences.lean}{antiDiagTensorEquiv_intertwines_rot,shiftSetCovariance_tensor_finrank}
Let $F$ be a field, $d \ge 1$, $A$ an $F$-vector space, and $p : \Z_d \to \kappa$.
\begin{enumerate}[label=(\alph*)]
\item A linear subspace $C \le \Mat_{d\times d}(A)$ satisfies \leanname{SharedDecoders.IsSharedCode}
  for $p$ if and only if for every shift $s \in \Z_d$ and every $Y \in C$ with $\rot_s(Y) \in C$,
  $Y$ lies in \lean{InPairwiseShiftCovariance}: whenever $p(c) = p(c+s)$, one has
  $Y_{c+s,\,j} = Y_{c,\,j+s}$ for all $j \in \Z_d$
  (cf.\ \leanname{SharedDecoders.isSharedCode_iff_compatible}).
\item Under $\Psi : (\Z_d \to F) \otimes_F (\Z_d \to A) \xrightarrow{\sim} \Mat_{d\times d}(A)$,
  column rotation $\rot_t$ intertwines with the pure right-factor coordinate shift
  $\mathrm{id} \otimes \sigma^t$ (\lean{coordShiftLinear}, $(\sigma^t v)_a = v_{a+t}$).
\item For any shift set $E \subseteq \Z_d$, a matrix $X \in \Mat_{d\times d}(A)$ lies in the
  shift-set covariance subspace $H_E(p)$ (\lean{shiftSetCovarianceSubmodule}, cf.\
  \leanname{SharedDecoders.covarianceSpace}) if and only if its
  anti-diagonal coordinates $Z_{c,a} = X_{c, a - c}$ are constant in $c$ on every connected component
  of \leanname{SwapRigidity.Link} $p$ $E$ (\lean{antiDiagonalTensorSeparation}). When
  $\dim_F A < \infty$, writing $U_E(p) \le (\Z_d \to F)$ for the subspace of vectors constant on
  those components (\lean{leftComponentSubmodule}), $\Psi$ maps $U_E(p) \otimes_F (\Z_d \to A)$ onto
  $H_E(p)$ and
  \[
    \dim_F H_E(p) \;=\; \dim_F U_E(p) \cdot d \cdot \dim_F A
  \]
  (cf.\ \leanname{SharedDecoders.finrank_covarianceSpace}).
\end{enumerate}
\end{theorem}
\begin{outline}
In~(a) (where linearity of $C$ is not used), determinism of $f_{p(c)}$ on $C$ means that whenever
$X_1, X_2 \in C$ satisfy $p(c_1) = p(c_2)$ and $\rot_{c_1}(X_1) = \rot_{c_2}(X_2)$, row $c_1$ of $X_1$
equals row $c_2$ of $X_2$: given $Y \in C$ with $\rot_s(Y) \in C$ and $p(c) = p(c+s)$, applying this
to $X_1 = \rot_s(Y) \in C$ at $c_1 = c$ and $X_2 = Y \in C$ at $c_2 = c + s$ (so
$\rot_c(X_1) = \rot_{c+s}(Y) = \rot_{c_2}(X_2)$) yields $(\rot_s Y)_{c, j} = Y_{c+s, j}$, i.e.\
$Y_{c, j+s} = Y_{c+s, j}$; conversely, if $\rot_{c_1}(X_1) = \rot_{c_2}(X_2)$ for $X_1, X_2 \in C$ with
$p(c_1) = p(c_2)$, setting $s = c_2 - c_1$ gives $X_1 = \rot_s(X_2)$, so taking $Y = X_2 \in C$ gives
$\rot_s(Y) = X_1 \in C$ and hence $(X_2)_{c_1+s, j} = (X_2)_{c_1, j+s} = (X_1)_{c_1, j}$. In~(b),
$\Psi(\mathrm{id} \otimes \sigma^t(u \otimes v))_{c,j} = u_c v_{j+c+t} = \rot_t(\Psi(u\otimes v))_{c,j}$.
In~(c), substituting $X_{c,j} = Z_{c, j+c}$ turns $X_{c', j+\delta'} = X_{c, j+\delta}$ under
$c + \delta = c' + \delta'$ into $Z_{c', a} = Z_{c, a}$ independently for each anti-diagonal
$a = c + \delta + j$.
\end{outline}

\begin{proposition}[Two-subspace Grassmann relaxation and alternating-mask shift separation; cf.\ {\cite[Lemma~3.1]{SharedDecoders2026}}]\label{prop:mask-floors}
\formal{DualityAndSwapConsequences.lean}{maskFloor_GenericSubspaceRealization}
Let $F$ be a field, $d \ge 1$, and $A$ a finite-dimensional $F$-vector space.
\begin{enumerate}[label=(\alph*)]
\item For any $p : \Z_d \to \kappa$, any shift $s \in \Z_d$, and any subspaces
  $C_0, C_s \le \Mat_{d\times d}(A)$ with $\dim_F C_0 = \dim_F C_s$ and
  $C_0 \cap C_s \le H_{\{0,s\}}(p)$, Grassmann modularity gives
  $2\dim_F C_0 \le d^2\dim_F A + \dim_F H_{\{0,s\}}(p)$. Conversely, once the right-factor shift
  coupling $C_s = \rot_{-s}(C_0)$ is dropped, choosing two independent subspaces of dimension
  $\lfloor(d - \dim_F U_{\{0,s\}}(p))/2\rfloor$ inside a complement of $U_{\{0,s\}}(p)$ in
  $(\Z_d \to F)$ constructs $C_0 = \Psi(V_0 \otimes \top)$ and
  $C_s = \Psi(V_s \otimes \top)$ with $C_0 \cap C_s = H_{\{0,s\}}(p)$, $\dim_F C_0 = \dim_F C_s$, and
  \[
    d^2\dim_F A + \dim_F H_{\{0,s\}}(p) \;\le\; 2\dim_F C_0 + d\dim_F A.
  \]
\item On even $d = 2m$ ($m \ge 2$) with the alternating parity partition
  $p : \Z_{2m} \to \mathrm{Bool}$ ($p(c) = \texttt{true} \iff c \text{ is odd}$, for
  $c \in \{0,\dots,2m-1\}$), shift
  $s = 1$ satisfies $p(c) \ne p(c+1)$ everywhere so $U_{\{0,1\}}(p) = \top$ and
  $H_{\{0,1\}}(p) = \top$, whereas shift $s = 2$ has $\dim_F U_{\{0,2\}}(p) = 2$ and forces
  $2\dim_F C \le ((2m)^2 + 4m)\dim_F A$ for every shared code $C$.
\end{enumerate}
\end{proposition}
\begin{outline}
Part~(a) applies $\dim(C_0 + C_s) + \dim(C_0 \cap C_s) = 2\dim C_0$ for the upper bound, and chooses
disjoint sets of $\lfloor(d - \dim U_{\{0,s\}}(p))/2\rfloor$ non-representative indices on the left
tensor factor for the converse construction. Part~(b) checks that shift~$1$ has no monochromatic
links on the alternating coloring while shift~$2$ connects all even indices and all odd indices into
two components.
\end{outline}
